\chapter{Block Definition Diagram}

\begin{lernziele}
Nach diesem Kapitel können Sie die Hauptstruktur des Roboters als Block Definition Diagram modellieren.
\end{lernziele}

\section{Zweck des BDD}
Das \gls{bdd} zeigt, aus welchen logischen oder physischen Bausteinen ein System besteht. Ein Block kann Hardware, Software, Mensch, Datenobjekt oder abstrakte Funktion darstellen.

\section{Hauptstruktur des Roboters}
\begin{verbatim}
Robot
|- NavigationSystem
|- PerceptionSystem
|- MotionControl
|- BatterySystem
|- CommunicationSystem
|- UserInterface
\end{verbatim}

Diese Darstellung ist noch grob. Sie eignet sich aber hervorragend, um den Systemumfang zu verstehen.

\section{Komposition}
Eine Komposition beschreibt eine starke Teil-Ganzes-Beziehung. Wenn das \texttt{PerceptionSystem} Teil des Roboters ist, kann dies als Komposition modelliert werden.

\section{Dekomposition}
Das \texttt{PerceptionSystem} kann weiter zerlegt werden:
\begin{verbatim}
PerceptionSystem
|- Camera
|- Lidar
|- ObjectDetector
|- ObstacleDetector
\end{verbatim}

\section{Logische und physische Blöcke}
Ein häufiger Anfängerfehler besteht darin, logische Funktionen und physische Bauteile unklar zu vermischen. \texttt{Camera} ist eher physisch. \texttt{ObjectDetector} ist eher logisch oder softwarebezogen. Beides darf im Modell vorkommen, sollte aber bewusst getrennt werden.

\begin{hinweisbox}
Ein BDD muss nicht sofort perfekt sein. Es entwickelt sich mit dem Verständnis des Systems.
\end{hinweisbox}

\begin{uebungbox}
Erstellen Sie ein BDD mit dem Block \texttt{Robot} und mindestens fünf Subsystemen. Ergänzen Sie beim \texttt{PerceptionSystem} Kamera, Lidar und Objekterkennung.
\end{uebungbox}
